%%=============================================================================
%% Conclusie
%%=============================================================================

\chapter{Conclusie}%
\label{ch:conclusie}

% TODO: Trek een duidelijke conclusie, in de vorm van een antwoord op de
% onderzoeksvra(a)g(en). Wat was jouw bijdrage aan het onderzoeksdomein en
% hoe biedt dit meerwaarde aan het vakgebied/doelgroep? 
% Reflecteer kritisch over het resultaat. In Engelse teksten wordt deze sectie
% ``Discussion'' genoemd. Had je deze uitkomst verwacht? Zijn er zaken die nog
% niet duidelijk zijn?
% Heeft het onderzoek geleid tot nieuwe vragen die uitnodigen tot verder 
%onderzoek?

In deze bachelorproef werd onderzocht op welke manier een geautomatiseerd systeem ontwikkeld kan worden voor het uitvoeren van onboarding- en offboardingprocessen binnen de hybride IT-infrastructuur van het lokaal gemeentebestuur van Berlare. Hierbij werd specifiek gekeken naar de automatisering van gebruikersbeheer binnen zowel Active Directory als Microsoft 365, Exchange Online en Entra ID.

Uit het onderzoek en de implementatie blijkt dat automatisering een duidelijke meerwaarde biedt binnen het beheer van gebruikersaccounts en toegangsrechten. Het ontwikkelde systeem slaagt erin om, vertrekkend van gegevens uit Microsoft Forms, automatisch gebruikersaccounts aan te maken, groepslidmaatschappen toe te kennen, licenties te beheren en mailboxrechten te configureren. Daarnaast werd eveneens een geautomatiseerd offboardingproces ontwikkeld waarbij gebruikersaccounts gedeactiveerd worden, rechten verwijderd worden en mailboxen automatisch omgezet worden naar gedeelde postvakken.

Door repetitieve beheertaken te automatiseren wordt de actieve verwerkingstijd voor de ICT-dienst aanzienlijk verminderd. Daarnaast zorgt de oplossing voor meer consistentie binnen de configuratie van gebruikersaccounts en wordt de kans op menselijke fouten beperkt. Hierdoor kunnen medewerkers van de ICT-dienst zich meer focussen op complexere technische taken en ondersteuning binnen de organisatie.

Naast de efficiëntiewinst draagt de oplossing ook bij aan een verbeterd beveiligings- en toegangsbeheer. Door gebruikers automatisch te verwijderen uit groepen, distributielijsten en mailboxrechten tijdens offboarding, wordt het risico verminderd dat voormalige medewerkers ongewenst toegang behouden tot interne systemen of gegevens.

Tijdens de ontwikkeling werd eveneens rekening gehouden met betrouwbaarheid en beheerbaarheid. Door gebruik te maken van Power Automate, Power Automate Desktop en herbruikbare PowerShell-scripts kon een centrale en relatief schaalbare oplossing opgebouwd worden die geïntegreerd kan worden binnen de bestaande hybride infrastructuur van de organisatie.\newline

Toch zijn er ook enkele beperkingen aanwezig binnen de huidige implementatie. Bepaalde configuraties, zoals specifieke rechten op netwerkschijven en hardwarevoorbereiding, verlopen nog steeds manueel. Daarnaast blijft de oplossing afhankelijk van synchronisatieprocessen tussen Active Directory, Entra ID en Exchange Online. Ook de huidige rechtenstructuur van bepaalde netwerkschijven binnen de organisatie vormt een uitdaging voor volledige automatisering, aangezien niet alle toegangen beheerd worden via veiligheidsgroepen.

Verder wordt binnen Exchange Online momenteel nog gebruikgemaakt van interactieve authenticatie. Hierdoor is dit onderdeel van de oplossing nog niet volledig automatisch. Een mogelijke uitbreiding voor de toekomst is het implementeren van certificaatgebaseerde authenticatie zodat ook deze configuraties volledig automatisch uitgevoerd kunnen worden. \newline

Ook heeft de gemeente twee omgevingen. Namelijk een omgeving voor de gemeente met e-mailadres ``@berlare.be'' en ook een voor het OCMW met als e-mail\-adres ``@ocmwberlare.be''. Hierdoor zijn er nog beperkingen voor het OCMW. De gehele bachelorproef is ontwikkeld voor de gemeente. De stromen kunnen hergebruikt worden om een systeem op te zetten voor het OCMW, maar dan moeten de PowerShell scripts aangepast worden. \newline

Samenvattend kan gesteld worden dat de onderzoeksvraag positief beantwoord werd. Het is mogelijk om een betrouwbaar en efficiënt geautomatiseerd systeem te ontwikkelen voor onboarding en offboarding binnen de hybride IT-omgeving van het lokaal gemeentebestuur van Berlare. De ontwikkelde oplossing vormt daarbij een sterke basis voor verdere uitbreiding en verdere optimalisatie van gebruikersbeheer binnen de organisatie.
